\subsection{Тестовое воздействие симметричным ЛЧМ сигналом с прямоугольной огибающей}

В заключительном пункте моделирования для CIC фильтра с корректором в качестве тестового воздействия применялся широкополосный симметричный линейно-частотно модулированный (ЛЧМ) сигнал с прямоугольной огибающей. Генерация базового видеосигнала осуществлялась с помощью встроенной функции \texttt{phased.LinearFMWaveform}. В соответствии с техническим заданием (Таблица 2) база сигнала была задана равной $B = 512$, а девиация частоты (полоса сигнала) согласована с рабочей полосой пропускания цифрового фильтра.

На рисунке \ref{fig:p6_lfm} представлены временные и спектральные характеристики зашумленного ЛЧМ-сигнала на промежуточной частоте (вход системы), а также результат его комплексной обработки — демодуляции, фильтрации и децимации (выход системы).  

\begin{figure}[H]
    \centering
    \includegraphics[width=0.9\textwidth]{p6_lfm_signal.png}
    \caption{Временные диаграммы и спектры ЛЧМ-сигнала до и после цифровой фильтрации (с учетом теплового шума)}
    \label{fig:p6_lfm}
\end{figure}

Для более детального анализа фазовой и спектральной структуры демодулированного сигнала дополнительно было проведено моделирование прохождения нормированного ЛЧМ-сигнала (без добавления аддитивного белого гауссова шума). Эти результаты визуализированы на рисунке \ref{fig:p6_video}.

\begin{figure}[H]
    \centering
    \includegraphics[width=0.9\textwidth]{p6_lfm_video_supplement.png}
    \caption{Временная структура и спектр нормированного ЛЧМ-видеосигнала в условиях без шума}
    \label{fig:p6_video}
\end{figure}

\textbf{Выводы по результатам моделирования:}
Временные диаграммы высокочастотного радиосигнала (\texttt{Signal LFM on IF}) демонстрируют заданную прямоугольную амплитудную огибающую, заполненную частотно-модулированным колебанием. После квадратурного переноса частоты и прохождения через систему децимирующих фильтров на выходе формируется комплексный видеосигнал (\texttt{Signal LFM on IF = 0} и \texttt{Signal LFM video}). 

На графике очищенного от шума видеосигнала (рис. \ref{fig:p6_video}, сверху слева) наглядно раскрывается внутренняя структура симметричного ЛЧМ-импульса: мгновенная частота колебаний в синфазном и квадратурном каналах (I и Q) линейно изменяется, плавно проходя через нулевое значение точно в середине временного окна (точка симметрии).

Спектральный анализ также подтверждает высокую эффективность спроектированной системы фильтрации. Исходный спектр ЛЧМ-сигнала, искаженный шумами АЦП (\texttt{Spectrum LFM on IF}), после обработки переносится на нулевую частоту (\texttt{Spectrum LFM on IF = 0}). На выходе отчетливо видна характерная для ЛЧМ-сигналов с большой базой плоская «прямоугольная» вершина спектра. Края этого спектра жестко и симметрично обрезаны амплитудно-частотной характеристикой многоскоростного фильтра, что свидетельствует об успешном подавлении внеполосного шума и зеркальных каналов при полном сохранении полезной широкополосной информации внутри полосы пропускания.